\documentclass[11pt]{article}
\usepackage[a4paper,margin=2.3cm]{geometry}
\usepackage[T1]{fontenc}
\usepackage[utf8]{inputenc}
\usepackage{lmodern}
\usepackage{amsmath,amssymb}
\usepackage{graphicx}
\usepackage{booktabs}
\usepackage{array}
\usepackage{xcolor}
\usepackage[authoryear,round]{natbib}
\usepackage[colorlinks=true,linkcolor=blue!50!black,citecolor=blue!50!black,urlcolor=blue!50!black]{hyperref}
\usepackage{caption}
\usepackage{setspace}
\usepackage{enumitem}
\captionsetup{font=small,labelfont=bf}
\setlist{itemsep=1pt,topsep=3pt}
\setstretch{1.08}

\newcommand{\todo}[1]{\textcolor{red}{[TODO: #1]}}
\input{tables/macros.tex}
\input{tables/descriptives_macros.tex}
\newcommand{\ratio}{\ensuremath{\rho}}

\title{Imputing diffusion profiles of white matter fiber tracts in early brain development:\\
\large an extensive evaluation of normative, local, graph and transformer models}
\author{Martin Styner$^{1}$, \todo{coauthors}$^{2}$\\[4pt]
\parbox{0.9\textwidth}{\centering\small $^{1}$Department of Psychiatry and Department of Computer Science, University
of North Carolina at Chapel Hill, NC, USA \todo{verify affiliation wording}\\ $^{2}$\todo{affiliations}}}
\date{\todo{date / corresponding author e-mail}}

\begin{document}
\maketitle

\begin{abstract}
\noindent Along-tract diffusion profiles are a common summary of white matter microstructure, but in infant and
toddler cohorts many profile values are missing: quality control rejects tracts that run through a regional
artifact, tracts are truncated at their ends, and a whole model fit (e.g.\ NODDI) can be unavailable for a scan.
Analyses that require complete data either drop these scans or fill the gaps ad hoc. We report an extensive
evaluation of imputation methods for such profiles, carried out while developing the \texttt{impute-profiles}
command of the open-source DTIPlayground package. On \bcpscans{} scans of \bcpsubjects{} children aged 0--65 months
from the Baby Connectome Project (\bcptracts{} atlas tracts, \bcppositions{} profile locations, FA, MD, RD, AD, NDI, ODI and FWF), we hid
measured values in five patterns (interior gaps, end gaps, whole profiles, all tracts through a spherical region,
and a whole family of measures) and compared several dozen settings of six method families on identical hidden values,
with 3 random draws of 54 validation scans and folds by subject. Imputation from a scan's own values beat the
age-bin mean only for gaps and missing ends along a profile; whole profiles and regions were 1.8 to 2.4 times as far
off. Imputing deviations
from the age-bin norm made every pattern better than the norm (overall RMSE \rPtwoSiren{} of the age-bin mean's). A
population model (probabilistic PCA of all tracts and measures) with a sine-activated network (SIREN) imputing its
residuals within 4\,mm reached \rDefault{}, and could impute a missing measure family (\rDefaultWdti{} for the tensor
measures, \rDefaultWnoddi{} for NODDI). Amortized networks that read a scan's observed values directly (a graph U-Net
and a transformer, both correcting the population prediction) were worse alone (\rUnet{}, \rTfFour{}) but much
better for a missing family (transformer \rTfFourWdti{} and \rTfFourWnoddi{}); blending them with the population
model plus SIREN, with weights per measure and distance to the nearest observed value, gave \rFinal{} (paired per-scan
ratio \cORatio{} against population model plus SIREN, better on \cOBetter\% of scan draws), and \rRefFinal{} when the
transformer was also trained on \hbcdscans{} scans of the HBCD study. Functa, a partial VAE, age-dependent PCA, acquisition
correction, context inputs and seed ensembles did not help, and pooling the second cohort into the population model
made it worse. Errors were lowest for MD, RD and NDI and highest for ODI, FWF and FA, and grew with the distance to the
nearest observed value. We give practical recommendations for imputing tract profiles in developmental cohorts.
\end{abstract}

\noindent\textbf{Keywords:} diffusion MRI; tractometry; fiber tract profiles; missing data; imputation; normative
modeling; infant brain development; transformer; graph neural network; probabilistic PCA

\section{Introduction}

Along-tract profiles sample diffusion measures such as fractional anisotropy (FA) or mean diffusivity (MD) at
regular arc-length positions along a fiber tract, and compare those profiles across subjects
\citep{goodlett2009group,colby2012alongtract,yeatman2012afq,verde2014unc}. They are widely used in studies of white matter development, where the
measures change fast and non-linearly over the first years of life \citep{geng2012tract,dubois2014early}, and they
increasingly include multi-compartment measures such as the neurite density (NDI), orientation dispersion (ODI) and
isotropic free water fraction (FWF) of NODDI \citep{zhang2012noddi}. Infant and toddler data are, however, incomplete:
motion, susceptibility and signal dropout corrupt local regions; quality control (QC) then rejects all tracts that
pass through such a region; tracts are truncated where the tractography or the atlas mapping misses their ends; and
some scans lack a model fit altogether. Many statistical tools for profiles (functional data analysis, mixed models
over locations, multivariate and machine-learning analyses of all tracts jointly,
\citealp{chamberland2019dimensionality,richiehalford2021afqinsight}) assume complete profiles, so that a scan
missing a few profile values is either discarded in full or patched by ad-hoc interpolation.

Imputation of missing values is a well studied problem \citep{rubin1976inference,little2019missing}, with general
methods ranging from nearest-neighbour \citep{troyanskaya2001knn} and chained-equation \citep{vanbuuren2011mice}
imputation to deep generative models \citep{mattei2019miwae,ma2019eddi}. Tract profiles have structure that generic
methods do not exploit: values vary smoothly along a tract; neighbouring tracts are dense in the brain; the
measures of one scan are strongly coupled; and, in a developmental cohort, the expected value at each location
depends strongly on age. It is not obvious which of these sources of information matters for which kind of missing
value, how much a normative reference adds, and whether modern neural architectures (implicit neural representations,
graph networks, transformers) bring anything beyond a well-specified linear population model.

This paper reports an extensive empirical evaluation of imputation methods for tract profiles in a normative infant
and toddler cohort, the Baby Connectome Project \citep[BCP;][]{howell2019bcp}. The evaluation was carried out in six
phases while developing the \texttt{impute-profiles} command of the open-source DTIPlayground package
\citep{dtiplayground}, and every phase set or confirmed the defaults of the released tool. The emphasis of this paper
is not methodological novelty but the knowledge the evaluation yields. We ask:
(i) which approaches work for which missingness patterns (interior gaps, end gaps, whole profiles, all tracts through a
region, a whole family of measures);
(ii) how much a normative model helps;
(iii) what population models, local three-dimensional (3-D) methods, graph networks and transformers each contribute;
(iv) what did not help (functa, a partial variational autoencoder, age-dependent population models, acquisition
correction, extra context inputs, seed ensembles);
(v) whether pooling a second, larger cohort (the HEALthy Brain and Child Development study, HBCD;
\citealp{volkow2021hbcd,dean2024hbcdmri}) helps, and how it should be done;
(vi) how the error varies by measure, age and pattern; and
(vii) what practitioners should do.
All methods were compared on identical hidden values, with three random draws of the validation scans, folds by
subject, and paired per-scan statistics.

\section{Related work}

\paragraph{Tract profiles.} Along-tract analysis was introduced for DTI by sampling diffusion tensor measures along
atlas or subject fiber bundles \citep{goodlett2009group,colby2012alongtract}, and popularized by automated pipelines
such as AFQ \citep{yeatman2012afq} and the UNC-Utah NA-MIC framework \citep{verde2014unc}, with tract segmentation by
learned models in recent pipelines \citep{wasserthal2018tractseg}. Tractometry has been extended to many measures at once
\citep{chamberland2019dimensionality}, and to multivariate and machine-learning analyses of all profiles of a scan
\citep{richiehalford2021afqinsight}. Its reliability has been studied \citep{kruper2021reliability}, but to our knowledge
the imputation of missing profile values has not been systematically evaluated.

\paragraph{Missing-data imputation.} Classical results distinguish missing completely at random, at random and not at
random \citep{rubin1976inference,little2019missing}. Practical imputation methods include $k$-nearest neighbours
\citep{troyanskaya2001knn}, multiple imputation by chained equations \citep{vanbuuren2011mice}, random forests
\citep{stekhoven2012missforest} and low-rank models; probabilistic principal component analysis (PPCA) handles missing
values in a principled way through the EM algorithm \citep{tipping1999ppca,ilin2010pca}. Deep generative models
amortize inference over incomplete data: the partial VAE of EDDI \citep{ma2019eddi} and MIWAE
\citep{mattei2019miwae}. Masked autoencoders \citep{he2022mae} learn to reconstruct hidden parts of an input with
transformers \citep{vaswani2017attention} and are, in effect, amortized imputers. For a broader review in pattern
recognition see \citet{garcialaencina2010pattern}.

\paragraph{Normative models.} Normative modeling characterizes the distribution of a brain measure in a reference
population as a function of covariates such as age, and expresses an individual as deviations from it
\citep{marquand2016normative,rutherford2022charting,bethlehem2022brainchart}. Our phase 2 imputes such deviations
instead of raw values; phase 3 adds a joint low-rank model of the deviations across all locations and measures.

\paragraph{Neural fields and functa.} Sine-activated networks (SIREN) represent signals as continuous functions of
coordinates \citep{sitzmann2020siren}. Functa \citep{dupont2022functa} and its medical variant MedFuncta
\citep{friedrich2025medfuncta} share one network across a dataset and fit a per-instance latent modulation, so that the
shared network acts as a prior; we test both a per-scan SIREN and a functa model.

\paragraph{Graph networks.} Graph convolutional networks \citep{kipf2017gcn} and graph U-Nets \citep{gao2019graphunet}
operate on irregular structures such as the locations of fiber tracts. CONFETTI \citep{confetti} builds a
multi-resolution neighbourhood graph over tract profile locations for scan classification; we adapt it to
imputation.

\paragraph{Harmonization.} Multi-site pooling of diffusion measures requires correcting site or acquisition effects,
classically by location-and-scale models such as ComBat \citep{johnson2007combat,fortin2017dti}; our batch correction
is an age-aware location-and-scale model of this kind.

\section{Materials and methods}

\subsection{Data}\label{sec:data}

\paragraph{Cohorts.} The normative cohort is the BCP \citep{howell2019bcp}: \bcpscans{} scans of \bcpsubjects{}
children, aged \bcpagemin{} to \bcpagemax{} months (median \bcpagemedian{}, quartiles \bcpageQOne{} and
\bcpageQThree{}; Table~\ref{tab:data}, Fig.~\ref{fig:data}a). \bcpsubjTwoPlus{} children contributed two or more scans
(up to \bcpsessmax{}). Scans fall in three age bins used throughout: 0--3 months (\bcpbinA{} scans), 4--9 months
(\bcpbinB{}) and 10--60 months (\bcpbinC{}). The scan names carry three acquisition labels
(\texttt{\bcpAcqAName}, \bcpAcqAN{} scans; \texttt{\bcpAcqBName}, \bcpAcqBN{}; \texttt{\bcpAcqCName}, \bcpAcqCN{}).
\todo{BCP acquisition details (scanner, b-values, directions, resolution), preprocessing pipeline, tensor and NODDI
fitting, ethics approval and consent statements.} The second cohort is the HBCD study
\citep{volkow2021hbcd,dean2024hbcdmri}: \hbcdscans{} scans of \hbcdsubjects{} children, median age \hbcdagemedian{}
months (range \hbcdagemin{} to \hbcdagemax{}; \hbcdbinA{}, \hbcdbinB{} and \hbcdbinC{} scans in the three bins, and
\hbcdbinnone{} outside them). \todo{HBCD release, age definition (the negative ages suggest age relative to a reference
date such as term-equivalent age), acquisition, processing and data use agreement.}

\paragraph{Profiles.} Profiles were extracted on the 69 tracts of the DTI\_IBISEP\_Feb26 atlas
\todo{atlas reference}, sampled every 1\,mm of arc length: 5,388 profile locations per measure. Each location has a
3-D position in atlas space, the point of the tract's axis (computed from the atlas fibers with
\texttt{dmrifiberprofile compute-axis}) at that arc length. Tract axes are dense: the median distance from a tract's
locations to the nearest other tract's axis is \axisTractMedian{}\,mm for the median tract, and \axisWithinTwo\% of all
locations are within 2\,mm (\axisWithinFour\% within 4\,mm) of another tract. Seven measures were analysed: FA, MD,
radial and axial diffusivity (RD, AD) of the diffusion tensor, and NDI, ODI and FWF of NODDI \citep{zhang2012noddi}
(FWF is the isotropic fraction of the NODDI fit). The free-water tensor measures
\citep{pasternak2009freewater}, imputed in the early phases, were left out after the free-water fits failed in many
HBCD scans (Section~\ref{sec:pooling}). The profiles were cleaned by the QC of DTIPlayground (\texttt{qc-profiles}:
outlier profiles removed, diffusivities above $3.6\times10^{-3}$\,mm$^2$/s removed as single values)
\todo{describe the QC criteria in more detail}.

\paragraph{Missingness of the real profiles.} After QC, \missCellShare\% of the BCP cells (\missCells{} of
\allCells{} scan $\times$ location $\times$ measure values) are missing. Nearly all
(\missWholeCellShare\%) belong to whole missing profiles: \missWholeN{} of \missProfilesN{} scan $\times$ tract
$\times$ measure profiles (\missWholeShare\%), always all measures of a tract at once (\missRejectedTracts{} rejected
scan $\times$ tract combinations). Only \missPartialN{} profiles are partly missing: \missPartialEnds{} miss their start
or end (median \missEndLenMedian{}, at most \missEndLenMax{} locations) and \missPartialInner{} have an interior gap.
\missScansAny{} of the \bcpscans{} scans miss at least one profile value; \missScansHalf{} scans miss more than half of
their tracts and are excluded by default. The share of scans with a given tract rejected ranges up to
\missTractMax\% (\missTractMaxName; median over tracts \missTractMedian\%; Fig.~\ref{fig:data}b). No BCP scan lacks
a whole measure family after QC, but this pattern is the norm in cohorts where one model fit fails or is not computed.

\begin{table}[t]
\centering\small
\caption{The two cohorts (computed from the profile tables).}\label{tab:data}
\input{tables/tab_data.tex}
\end{table}

\begin{figure}[t]
\centering
\includegraphics[width=\textwidth]{figures/fig_data.pdf}
\caption{Data. (a) Ages at scan of BCP (normative cohort) and HBCD (reference cohort); dashed lines: the age bins.
(b) Real missingness of BCP after QC: share of scans with a tract rejected (all measures missing), against the
median distance from the tract's axis to the nearest other tract.}\label{fig:data}
\end{figure}

\subsection{Notation and the age-bin norm}

For scan $s$, location $p$ (tract, arc length) and measure $m$, $y_{spm}$ is the profile value. The age-bin norm is the
mean $\mu_{bpm}$ and SD $\sigma_{bpm}$ of the cohort's scans of age bin $b$; the deviation of a scan is
$z_{spm}=(y_{spm}-\mu_{b(s)pm})/\sigma_{b(s)pm}$. Methods ``on the deviations'' impute $z$ and return
$\hat y=\mu+\sigma\hat z$ with the norm of the missing location. In validation, the norm is always computed
without every scan of the validated scan's subject.

\subsection{Imputation methods}\label{sec:methods}

Table~\ref{tab:methods} summarizes the methods, grouped by the six phases of the evaluation. All are implemented in
\texttt{dtiplayground/dmri/fiberprofile/analysis/impute\_profiles.py} (with \texttt{graph\_unet.py},
\texttt{amortized.py} and \texttt{functa.py}).

\begin{table}[t]
\centering\small
\caption{Methods and what they use. Patterns: G interior gap, E end gap, W whole profile, R region, F whole measure
family. ``Own'': values of the scan itself; ``norm'': age-bin norm; ``cohort'': models fitted on the normative
cohort (without the validated subject).}\label{tab:methods}
\begin{tabular}{@{}p{3.3cm}p{2.0cm}p{7.6cm}l@{}}
\toprule
Method & Information & Description & Patterns \\
\midrule
Arc-linear & own & linear interpolation along the arc length of the same profile; ends left, repeated (constant)
or extrapolated & G, E \\
Linear (3-D) & own & barycentric interpolation in the Delaunay tetrahedra of the scan's observed locations of the
same measure in all tracts & G, E, W, R \\
$k$NN & own & inverse-distance-weighted mean of the $k$ nearest observed locations in 3-D & G, E, W, R \\
SIREN & own & per-scan sine network, inputs 3-D position (xyz), plus arc length or one-hot tract; all measures jointly
& G, E, W, R \\
Phase 2: any of the above on $z$ & own + norm & imputes deviations from the age-bin norm & as above \\
Population model (PPCA) & own + cohort & probabilistic PCA of the deviations of all locations and measures jointly;
posterior mean and SD of the missing cells & all \\
PPCA + local & own + cohort & SIREN or $k$NN imputes the residuals from the PPCA prediction within 4\,mm of an observed
value; PPCA beyond & all \\
Graph U-Net (hybrid) & own + cohort & multi-resolution graph convolutional U-Net over the locations; corrects the PPCA
prediction & all \\
Transformer (hybrid) & own + cohort & masked autoencoder over patches of 8 locations; corrects the PPCA prediction &
all \\
Blend & own + cohort & $(1-w)\,\hat y_{\mathrm{PPCA+SIREN}}+w\,\hat y_{\mathrm{model}}$, $w$ fixed, by precision, or per
measure and distance band & all \\
\midrule
\multicolumn{4}{@{}l}{\textit{Explored, not adopted}} \\
Functa & own + cohort & shared SIREN, per-scan latent fitted by meta-learned gradient steps & all \\
Partial VAE & own + cohort & set encoder of the observed cells, per-location decoder (EDDI/MIWAE-style) & all \\
Age-dependent PPCA & own + cohort & PPCA per age, scans weighted by a kernel of log age & all \\
Acquisition correction & own + cohort & random-intercept offset per acquisition label, removed before imputation &
all \\
\bottomrule
\end{tabular}
\end{table}

\paragraph{Phase 1: the scan's own values.} \emph{Arc-linear} interpolates linearly along the arc length between the
observed values at both ends of a gap (gaps up to 15\,mm); missing ends are left, filled with the nearest observed
value (\emph{constant}) or extrapolated from the three nearest. \emph{Linear} interpolates barycentrically within the
Delaunay tetrahedra of the observed 3-D positions of the same measure in all tracts. \emph{$k$NN} averages the $k$
nearest observed positions with inverse-distance weights. \emph{SIREN} \citep{sitzmann2020siren} fits, per scan, a sine
network (3 hidden layers of 256 units, Adam, learning rate $10^{-4}$; \citealp{kingma2015adam}) to all observed values
of all measures jointly, each standardized within the scan; its inputs are the 3-D position scaled to $[-1,1]$
(\emph{xyz}), optionally with the arc length (\emph{xyz+arc}) or a one-hot tract code (\emph{xyz+tract}); we varied
$\omega_0\in\{5,10,20,30\}$ and 100--1000 epochs.

\paragraph{Phase 2: deviations from the age-bin norm.} The same methods applied to $z$ instead of $y$. Values farther
than a distance limit from any observed value of the same measure (20\,mm on the deviations) are left to the age-bin
mean.

\paragraph{Phase 3: population model.} The deviations of the normative scans, over all locations and measures jointly
($5{,}388\times7$ cells observed in at least 10 scans), are modelled by PPCA \citep{tipping1999ppca},
$\mathbf z=\boldsymbol\mu+W\mathbf x+\boldsymbol\epsilon$, $\mathbf x\sim\mathcal N(0,I_q)$,
$\boldsymbol\epsilon\sim\mathcal N(0,\psi I)$, fitted by EM on the observed cells with the exact posterior of every row
\citep{ilin2010pca}, a ridge prior on the rows of $W$, and $q=80$ components by default (5--160 tested; isotropic,
per-measure and per-cell noise tested). A scan's observed deviations give the posterior mean and SD of every missing
cell, including all cells of a measure family the scan lacks. The posterior SDs are rescaled per measure so that
errors on held-out whole profiles (5 folds by subject) have unit root-mean-square $z$. The \emph{PPCA + local}
combination uses the PPCA prediction as the scan's own norm and imputes the residuals with SIREN (xyz, $\omega_0=10$,
500 epochs) or $k$NN ($k=20$) within 4\,mm of an observed value; beyond, the PPCA prediction is used. This combination
(PPCA + SIREN) was the default of DTIPlayground up to version 0.9.5.

\paragraph{Phase 3b: continuous age norms and pooling.} Instead of age bins, the mean and SD of every cell can be smooth
functions of $\log(\text{age}+1\,\text{month})$: a natural cubic spline or a local-linear kernel regression with
leave-one-out bandwidth. To pool cohorts, each batch is mapped to the reference batch by an age-aware shift and scale
per cell and measure (one spline age curve for all batches plus a shift per batch, fitted jointly; the scale as the
ratio of residual SDs over the common age range), and imputed values are mapped back.

\paragraph{Phase 4: graph U-Net.} After CONFETTI \citep{confetti}, a multi-resolution graph keeps every $2^L$-th
location of each tract at level $L$ (4 levels); nodes are linked along the tract and to the closest node of each
other tract within 1.5 times the median distance to the nearest other tract's axis, with Gaussian edge weights. A
U-Net \citep{gao2019graphunet} of graph convolutions \citep{kipf2017gcn} (128 features, one convolution per level,
batch normalization, dropout 0.1, skip connections) receives per location the observed deviations (0 where missing),
the observation mask, the 3-D position and arc length, and, in the \emph{hybrid} variant, the PPCA prediction and its
log SD; a scan-wide summary (mean and max at the coarsest level) is added there. It outputs a mean and log variance per
location and measure and is trained (AdamW, \citealp{loshchilov2019adamw}; cosine schedule; batches of 16 scans; up to
60 epochs) with the validation patterns of Section~\ref{sec:protocol} hidden anew in every epoch (a region in half of
the scans, a whole family in a quarter), on the Gaussian negative log-likelihood of the hidden values
\citep{kendall2017uncertainties} (visible values with weight 0.1). Ten percent of the training subjects select the
epoch and calibrate the SDs. The PPCA input of a training scan comes from a model fitted without its subject (5-fold
cross-fitting). We also evaluated CONFETTI's original configuration (graph threshold from the dense atlas fibers, 64
features, no scan-wide summary).

\paragraph{Phase 6: transformer.} A masked-autoencoder transformer \citep{vaswani2017attention,he2022mae}: each tract
is cut into patches of 8 consecutive locations (about 700 tokens per scan); a token embeds the per-location inputs of
its patch (as for the U-Net) plus a learned tract embedding; 6 pre-norm self-attention layers of width 192 (6 heads,
dropout 0.1) relate every token to every other, and a linear head returns a mean and log variance for every location of
the patch. Training is as for the graph U-Net (AdamW, learning rate $10^{-3}$, weight decay $10^{-4}$, cosine schedule,
200 or 400 epochs with the best held-out epoch kept), in the hybrid variant correcting a cross-fitted PPCA prediction.

\paragraph{Phase 6: blends.} The PPCA + SIREN prediction $b$ and a model prediction $g$ (U-Net or transformer) are
combined as $(1-w)b+wg$, with (i) fixed weights $w_{\text{near}}$ for cells within 4\,mm (those SIREN fills) and
$w_{\text{beyond}}$ for the others; (ii) precision weights $w=\sigma_b^{-2}/(\sigma_b^{-2}+\sigma_g^{-2})$; or (iii) a
\emph{weight table} per measure $m$ and band $k$ of the distance to the nearest observed value of the measure (0--0.5,
0.5--1, 1--2, 2--3, 3--4\,mm for SIREN cells; 4--6, 6--10, 10--20, $>$20\,mm for fallback cells; and a measure the scan
has no observed value of). The table is the weighted least-squares solution in deviation units,
\begin{equation}
w_{mk}=\operatorname{clip}_{[0,1]}\frac{\sum_{i\in(m,k)}c_i\,(t_i-b_i)(g_i-b_i)}{\sum_{i\in(m,k)}c_i\,(g_i-b_i)^2},
\end{equation}
over hidden cells $i$ with true deviation $t_i$, where $c_i$ is the inverse mean squared error of the age-bin mean
in the cell's (pattern, measure, age bin) group, so that each group counts as in the score. This is a constrained form
of stacking \citep{wolpert1992stacked,breiman1996stacked}. The table shipped with DTIPlayground 0.9.8 was fitted on
all validation cells of the three draws; all scores reported for it here are cross-fitted (weights for a scan fitted
without its subject, 5 folds). We also tested a gradient-boosted weight per cell and calibrated precision weights.

\paragraph{Reference cohort.} With \texttt{--transformer-reference-dir}, a second cohort (HBCD) is batch-corrected to
BCP at the same age and added to the training of the transformer only; the norm, PPCA and SIREN remain those of BCP.
We compared this with pooling BCP and HBCD for the norm, the PPCA and the transformer.

\paragraph{Explored alternatives.} \emph{Functa} \citep{dupont2022functa,friedrich2025medfuncta}: one SIREN shared
by the normative scans, modulated per scan by shifts from a latent vector fitted by a few meta-learned gradient steps
on the scan's observed values, meta-trained on the hidden values of the validation patterns; we varied inner steps and
step sizes, test-time refinement, an encoder initialization, depth (2--8), width (64--512), latent size (16--512), age
and tract inputs, and graph variants (Laplacian-eigenvector inputs, a learned vector per location, a graph-convolution
decoder), each on the deviations and as a hybrid correcting PPCA. \emph{Partial VAE}
\citep{ma2019eddi,mattei2019miwae}: a set encoder of the observed cells (per-location MLP pooled by mean and max), a
latent (32--128), a decoder with weights per location, importance-weighted training. \emph{Age-dependent PPCA}: a PPCA
per age, the scans weighted by a Gaussian kernel of log age. \emph{Acquisition correction}: an offset per acquisition
label and cell fitted on the deviations with a random intercept per subject, removed before imputation. \emph{Context
inputs} to the transformer: distance to the nearest observed value, age and acquisition. \emph{Seed ensembles}: three
transformers averaged \citep{lakshminarayanan2017ensembles}.

\subsection{Evaluation protocol}\label{sec:protocol}

\paragraph{Hiding patterns.} Measured values are hidden, imputed, and compared with the truth (Fig.~\ref{fig:patterns}).
Each validated scan goes through three rounds; each round has three passes on the complete scan:
(1) whole profiles (all measures) of 10\% of the tracts with at least three observed locations; of the remaining tracts,
an interior gap of 3--15\,mm in 20\% (\emph{gap}) and the first or last 1--8\,mm of the observed part in another 20\%
(\emph{end gap}; at least three observed locations kept); (2) every tract with an observed location inside a sphere of
radius 10--25\,mm around a random observed location loses its whole profile (\emph{region}; spheres reaching more than
half of the tracts are redrawn), mimicking QC rejection after a regional artifact; (3) all values of one measure family
the scan has, the tensor family (FA, MD, RD, AD) or the NODDI family (NDI, ODI, FWF), chosen at random (\emph{whole
family}). The hiding is identical for all settings of a run and across runs with the same random seed; whole profiles
and regions mirror the dominant real missingness (Section~\ref{sec:data}), end gaps its most common partial pattern.

\begin{figure}[t]
\centering
\includegraphics[width=\textwidth]{figures/fig_patterns.pdf}
\caption{Missingness patterns and evaluation protocol. (a, b) One round of hidden values on a 12-month BCP scan, the
profile locations projected on axial and sagittal planes (grey: observed). (a) Pass 1 hides whole profiles of random
tracts, interior gaps and end gaps; (b) pass 2 hides all tracts through a sphere (this draw: 13 of 69 tracts). Pass 3
(not shown) hides a whole measure family. (c) The protocol.}\label{fig:patterns}
\end{figure}

\paragraph{Sampling and folds.} Per random seed (0, 1, 2), 20 scans per age bin were drawn, at most half of a bin
(\nScansBinA{}, \nScansBinB{} and \nScansBinC{}: \nScansSeed{} scans), \nScansDistinct{} distinct scans of
\nSubjDistinct{} subjects over the three seeds, \nPairs{} scan $\times$ seed pairs. Every model fitted on the cohort
(norm, PPCA, continuous norms, batch correction, networks, blend weights) is fitted without the validated subjects:
the validated scans are split by subject into 5 folds and each fold's models are fitted on all other subjects.

\paragraph{Score.} For a setting, pattern $P$, measure $m$ and age bin $b$, the RMSE over the hidden values is divided by
the RMSE of the age-bin mean on the same values,
\begin{equation}
\ratio_{Pmb}=\mathrm{RMSE}_{Pmb}(\text{setting})/\mathrm{RMSE}_{Pmb}(\text{age-bin mean}),
\end{equation}
so that $\ratio<1$ means better than the norm. Values a setting leaves unfilled (e.g.\ whole profiles for arc-linear,
cells beyond a distance limit, a whole family for same-scan methods) are counted with the age-bin mean, as an imputation
with that fallback would do. The pattern score is the mean of $\ratio$ over measures and age bins; the \emph{overall}
score is the mean over the four patterns gap, end gap, whole profile and region; the whole-family pattern is reported
separately for the tensor and the NODDI family. Scores are averaged over the three seeds (SD between seeds in
parentheses). Absolute RMSEs are reported in the units of each measure.

\paragraph{Paired statistics.} Seed-to-seed variation of a score (SD up to about 0.01 overall and more within an age
bin) is as large as many differences between settings, so settings are compared on the same scans. The per-scan score
is the mean over pattern $\times$ measure of $\sqrt{\mathrm{MSE}/\mathrm{MSE}_{\text{age-bin mean}}}$ for that scan; for
two settings A (reference) and B, we report the geometric mean over scan $\times$ seed pairs of the ratio of their
per-scan scores, its 95\% confidence interval from the normal approximation of the log ratios, the Wilcoxon signed-rank
test \citep{wilcoxon1945} of the log ratios, and the share of pairs in which B is better. The \nPairs{} pairs are not
independent: the same scan can be drawn by several seeds, and several scans come from one subject. We therefore also
report a cluster bootstrap CI that resamples the \nSubjPairs{} subjects (4,000 resamples; \citealp{field2007bootstrap})
and a Wilcoxon test on per-subject mean log ratios, and we Holm-adjust \citep{holm1979} the $p$ values over the
\nContrasts{} pre-specified contrasts of Table~\ref{tab:paired}. Paired tests are only reported between settings that
fill every hidden value.

\paragraph{Quick screens.} Model families that failed to justify a full validation (functa, partial VAE,
age-dependent PPCA) and transformer variants were screened on one fixed held-out split: 30 test subjects, two rounds
of the same hiding patterns, models fitted on the other subjects, and the same score on the 7 measures. Screen numbers
are comparable within the screen but not directly with the 3-seed validation.

\paragraph{Earlier protocols.} Phases 1--3 were first run on 30 scans, one seed and 11 measures (including the
free-water tensor measures), and phase 3b on the 11 measures with three seeds. All phases 1--4 were later re-run with
the protocol above; we report those results throughout, and use the earlier runs only where no re-run exists (the
SIREN hyperparameter grid, the distance limit of the PPCA residuals and the pooling experiment; Appendix).

\subsection{Implementation and reproducibility}

All methods are part of DTIPlayground 0.9.8 \citep{dtiplayground} (PyTorch, \citealp{paszke2019pytorch}). Each
validation run is a single \texttt{dmrifiberprofile impute-profiles --validate} call, with \texttt{--validate-seed}
selecting the sample and the hidden values; its outputs include per-scan, per-age-bin and per-distance error tables
and, with \texttt{--validate-dump-}\allowbreak\texttt{dir}, every hidden cell with all predictions, from which the blends were fitted
offline. A run with the population model, SIREN and one network took about 2 hours on one NVIDIA RTX A6000 GPU.
Every number, table and figure of this paper is produced from those result files by the scripts accompanying the
paper (Data and code availability).

\section{Results}

\subsection{Overview: from the scan's own values to blended models}

Table~\ref{tab:main} and Fig.~\ref{fig:progression} give the error of representative settings of every phase on the
same hidden values.

\begin{table}[p]
\centering\footnotesize
\caption{Error relative to the age-bin mean (RMSE ratio $\ratio$; below 1: better) on identical hidden values,
mean over 3 seeds $\times$ 54 BCP scans, 7 measures and 3 age bins. Overall: mean of the four patterns (SD between
seeds). $^{\dagger}$Leaves some hidden values, counted with the age-bin mean (a whole family is never filled by
phases 1--2). $^{\ast}$Weights cross-fitted by subject; the mean of the per-seed scores (pooled over seeds, as in the
software documentation: 0.567 and 0.538).}\label{tab:main}
\setlength{\tabcolsep}{4pt}
\resizebox{\textwidth}{!}{\input{tables/tab_main.tex}}
\end{table}

\begin{figure}[t]
\centering
\includegraphics[width=\textwidth]{figures/fig_progression.pdf}
\caption{Phase-wise progression. (a) Error relative to the age-bin mean per pattern (markers) and overall (bars). (b)
A whole measure family hidden; methods of phases 1--2 cannot impute it (1.0 by construction). P1--P6: phases;
``released'': the default of DTIPlayground 0.9.8 (scored cross-fitted).}\label{fig:progression}
\end{figure}

\paragraph{Phase 1: the scan alone is not enough.} Imputed from the scan's own values, only interior gaps were better
than the age-bin mean (arc-linear \rPoneArcGap{}, SIREN xyz \rPoneSirenGap{}); with constant ends, arc-linear also
improved end gaps (\rPoneArcEnd{}). Whole profiles and regions were imputed far worse than the norm by every 3-D
method (SIREN xyz \rPoneSirenWhole{} and \rPoneSirenRegion{}, linear \rPoneLinearWhole{} and \rPoneLinearRegion{}).
Neighbouring tracts differ systematically in their measures, by orientation and organization, so a missing tract
cannot be read off its neighbours' values. The best phase-1 setting overall, arc-linear with constant ends
(\rPoneArc{}), owes this to filling nothing beyond gaps and ends.

\paragraph{Phase 2: the norm makes every pattern imputable.} On the deviations from the age-bin norm, every 3-D method
was better than the age-bin mean for whole profiles, and all but $k$NN with $k=1$ for regions. The best single setting was SIREN
with the 3-D position only (\rPtwoSiren{} overall; gaps \rPtwoSirenGap{}, end gaps \rPtwoSirenEnd{}, whole profiles
\rPtwoSirenWhole{}, regions \rPtwoSirenRegion{}). A scan that deviates from the norm in one tract tends to deviate
the same way in its neighbours, while the norm carries the tract-specific level that the neighbours cannot. The gain
shrank as more was missing: regions remained close to the norm.

\paragraph{Phase 3: a population model.} PPCA alone (\rPpca{}) was better than phase-2 SIREN for whole profiles
(\rPpcaWhole{} vs \rPtwoSirenWhole) and regions (\rPpcaRegion{} vs \rPtwoSirenRegion), but worse for gaps and ends,
where neighbours along the same profile carry more. Imputing the PPCA residuals locally combined both: PPCA + SIREN
(\rDefault{}; gaps \rDefaultGap{}, end gaps \rDefaultEnd{}, whole profiles \rDefaultWhole{}, regions
\rDefaultRegion{}) was better than PPCA alone on every scan draw (paired ratio \cARatio{}, Table~\ref{tab:paired}) and
better than PPCA + $k$NN (\rPpcaKnn{}; $k$NN was \cBRatio{} times SIREN's error, although slightly better for whole
profiles and regions). PPCA is also the first method that imputes a missing measure family (tensor \rPpcaWdti{},
NODDI \rPpcaWnoddi{}).

\paragraph{Phase 4: graph U-Net.} In CONFETTI's original configuration the graph U-Net imputed poorly (0.90 of the
age-bin mean's error overall in a development run on one seed and 11 measures, against 0.69 for PPCA alone; software
documentation). Three changes made it competitive: a graph built from the tract axes rather than the
dense fibers (fewer edges to other tracts), a scan-wide context at the coarsest level, and correcting the PPCA
prediction (hybrid). The hybrid U-Net alone (\rUnet{}) was better than PPCA alone for every pattern but worse than
PPCA + SIREN for gaps and ends (\rUnetGap{} vs \rDefaultGap{}); adding SIREN on its residuals did not close the gap
(\rUnetSiren{}). Its strength was a missing family: \rUnetWdti{} for the tensor and \rUnetWnoddi{} for the NODDI
family, against \rPpcaWdti{} and \rPpcaWnoddi{} for PPCA. Used only beyond 4\,mm (PPCA + SIREN near, U-Net beyond)
it changed the four patterns little (\rUnetBeyond{}; paired \cFRatio{}).

\paragraph{Phase 6: transformer and blends.} The transformer alone (\rTfTwo{} with 200 epochs, \rTfFour{} with 400)
was better than the U-Net alone (paired \cIRatio{}) and better still for a missing family (tensor \rTfFourWdti{},
NODDI \rTfFourWnoddi{} with 400 epochs). Neither network alone reached PPCA + SIREN on the four patterns (transformer
\cHRatio{}, U-Net \cERatio{} times its error), because SIREN on the PPCA residuals is better close to observed values.
Blending did: a fixed blend with weight 0.25 near and 0.75 beyond 4\,mm reduced the error to \rUnetBlend{} with the
U-Net and \rTfBlend{} with the transformer (\rTfBlendFour{} with 400 epochs), and the learned weight table to
\rFinal{} (cross-fitted). Against PPCA + SIREN, the released default (transformer of 400 epochs, weight table) had a
paired ratio of \cORatio{} and was better on \cOBetter\% of scan draws; it improved every pattern (gaps
\rFinalGap{}, end gaps \rFinalEnd{}, whole profiles \rFinalWhole{}, regions \rFinalRegion{}) and a missing family by
far (tensor \rFinalWdti{}, NODDI \rFinalWnoddi{}).

\subsection{Which approach works for which pattern}

\paragraph{Interior gaps and end gaps.} Gaps are the easiest pattern: \shareFourGap\% of the hidden values are within 4\,mm of an observed
value (\shareFourEnd\% for end gaps, \shareFourWhole\% for whole profiles, \shareFourRegion\% for regions;
Fig.~\ref{fig:distance}d), and local methods on deviations or PPCA residuals win. Without a norm, arc-linear
interpolation along the profile is a reasonable choice for interior gaps (\rPoneArcGap{}); with a norm, arc-linear
on the deviations reached \rPtwoArcGap{} and SIREN on the PPCA residuals \rDefaultGap{}. End gaps behave like gaps
with one-sided support and were consistently about 0.1 worse. Amortized networks alone were clearly worse here
(transformer \rTfFourGap{}), but contributed in the blend (\rFinalGap{}).

\paragraph{Whole profiles and regions.} A whole profile is imputed from neighbouring tracts and from the scan's other
tracts and measures; a region also removes the neighbours. Here the population model and the amortized networks
mattered: PPCA alone \rPpcaWhole{} and \rPpcaRegion{}, phase-2 SIREN \rPtwoSirenWhole{} and \rPtwoSirenRegion{}, the
transformer alone \rTfFourWhole{} and \rTfFourRegion{}, the released blend \rFinalWhole{} and \rFinalRegion{}. Regions
remain the hardest of the four patterns for every method.

\paragraph{A whole measure family.} Only cohort models can impute a measure the scan lacks entirely, from the
scan's other measures. The amortized networks were much better than PPCA (tensor family: transformer \rTfFourWdti{},
U-Net \rUnetWdti{}, PPCA \rPpcaWdti{}; NODDI family: \rTfFourWnoddi{}, \rUnetWnoddi{}, \rPpcaWnoddi{}); paired over
scans, the transformer's error for a missing family was \wCRatio{} of PPCA's (400 epochs, better on \wCBetter\% of scan
draws). The tensor family was easier to recover from NODDI than the reverse, and MD and RD were recovered best (per
measure: Table~\ref{tab:metric}).

\paragraph{Distance to the nearest observed value.} Within each pattern, the error grew with the distance to the nearest
observed value of the same measure (Fig.~\ref{fig:distance}). The deviations of a far neighbour still carried
information on the deviations scale (phase-2 SIREN stayed below 1 up to its 20\,mm limit for regions), but the PPCA residuals
are local: SIREN on them helped within a few millimetres and hurt beyond, which is why it is limited to 4\,mm
(Appendix Fig.~\ref{fig:cap}). Beyond 4\,mm the amortized networks were better than PPCA, which is what the blend
exploits (Fig.~\ref{fig:weights}c).

\begin{figure}[t]
\centering
\includegraphics[width=\textwidth]{figures/fig_distance.pdf}
\caption{Error by the distance from a hidden value to the nearest observed value of the same measure, per pattern
(3 seeds; bins with fewer than 600 hidden values omitted; vertical line: 4\,mm). (d) The distribution of these distances
per pattern. Interior gaps are omitted from (a)--(c): \shareFourGap\% of their values lie within 4\,mm.}\label{fig:distance}
\end{figure}

\subsection{How much does the normative model help?}

The normative information enters at three levels. The age-bin norm alone turned SIREN's whole profiles from \rPoneSirenWhole{} (phase 1, on
values) into \rPtwoSirenWhole{} (phase 2, on deviations), its regions from \rPoneSirenRegion{} into
\rPtwoSirenRegion{}, and its overall score from \rPoneSiren{} to \rPtwoSiren{}. The population model, which captures how deviations co-vary across all tracts and
measures, took the overall error from \rPtwoSiren{} to \rDefault{} and made a missing family imputable. Learned
amortized models of the deviations took it to \rFinal{} (or \rRefFinal{} with a larger training cohort). In absolute
terms (Table~\ref{tab:absolute}), the released default roughly halved the RMSE of the age-bin mean for the
diffusivities.

A continuous age norm (spline or kernel regression of log age) was a better predictor than the age-bin mean on its own
(paired \vpKernNormRatio{}, 95\% CI \vpKernNormCi{}; score \vNormKernelA{} in 0--3 months, where the norm
changes fastest within a bin; Appendix Table~\ref{tab:pooling}), but under PPCA + SIREN it made no difference (\vpKernRatio{}, CI \vpKernCi{}): the
population model and the scan's own values absorb the within-bin age trend.

\subsection{Population model components and local methods}

The PPCA error fell up to 80 components and was flat beyond (overall alone: \rPpcaTen{} with 10, \rPpcaForty{} with
40, \rPpca{} with 80, \rPpcaTwelve{} with 120). Combined with SIREN, 40 and 120 components were slightly but
consistently worse than 80 (paired \cCRatio{} and \cDRatio{}). The three noise models differed by less than 0.01 in
the earlier 30-scan runs (software documentation), so the isotropic model was kept. Within SIREN, the trade-off between
patterns was clear in the hyperparameter grid (Appendix Fig.~\ref{fig:sirengrid}): gaps and ends need long training and
a high $\omega_0$ or a tract code, whereas whole profiles and regions are best after 100 epochs and get worse as the
network fits the observed tracts more closely; xyz with $\omega_0=10$ and 500 epochs balanced them best on the
deviations.

\subsection{Blend weights}

The learned weight of the transformer rose with the distance to the nearest observed value
(Fig.~\ref{fig:weights}): \wtCNearMin--\wtCNearMax{} next to an observed value (where SIREN on the PPCA residuals
dominates), about 0.6 from 3--4\,mm, \wtCFarMin--\wtCFarMax{} for fallback cells, and \wtCFamMin--\wtCFamMax{} for a
missing family (BCP-trained transformer). The transformer trained with HBCD received more weight next to observed values
(\wtRNearMin--\wtRNearMax). Refitting the weight table from the validation dumps reproduced the table shipped in
DTIPlayground 0.9.8 to two decimals (largest difference \weightTableMaxDiff). Cross-fitted, the table was better than
the fixed 0.25/0.75 weights with the same transformer (paired \cNRatio{}, better on \cNBetter\%), and 400 training
epochs were better than 200 (\cMRatio{}). Precision weights from the predicted SDs were worse than fixed weights
(\rPrecisionBlend{} vs \rTfBlend{}; paired \cLRatio{}), because the SD attached to a SIREN value is that of the
population model, an upper bound next to an observed value; a gradient-boosted weight per cell and calibrated
precision weights were about as good as the table (software documentation). Among fixed weights,
0.25/0.75 (\rTfBlend{}) and 0.25/0.5 (\rTfBlendQuarterHalf{}) were best; weight 0 near observed values was clearly worse
(\rTfBlendZeroThree{}).

\begin{figure}[t]
\centering
\includegraphics[width=\textwidth]{figures/fig_weights.pdf}
\caption{Blend weights. (a, b) Weight $w$ of the transformer per measure and distance band (SIREN cells within
4\,mm, fallback cells beyond, ``none'': the scan has no observed value of the measure), fitted on all validation cells;
(a) transformer trained on BCP (the default table), (b) trained on BCP and HBCD. (c) Error relative to the age-bin
mean against distance, four patterns pooled, for PPCA + SIREN, the BCP transformer alone and the cross-fitted table
blends.}\label{fig:weights}
\end{figure}

\subsection{What did not help}

\begin{table}[t]
\centering\footnotesize
\caption{Quick held-out screen (one split, 30 test subjects, 2 rounds; 7 measures): best result per model family and
target; ``Runs'': number of configurations and test-time variants screened. Not directly comparable with
Table~\ref{tab:main}.}\label{tab:screens}
\setlength{\tabcolsep}{4pt}
\input{tables/tab_screens.tex}
\end{table}

\begin{figure}[t]
\centering
\includegraphics[width=0.75\textwidth]{figures/fig_screens.pdf}
\caption{Quick held-out screen: overall error of every screened configuration (dots; bar: best) of the model
families that were not adopted and of the transformer, with PPCA alone and the hybrid graph U-Net on the same
split.}\label{fig:screens}
\end{figure}

\paragraph{Functa.} Over \nFunctaConfigs{} screened runs (configurations and test-time variants), the best functa
model reached \sFunctaHyb{} as a hybrid correcting PPCA and \sFunctaDev{} on the deviations, against \sPpca{} for PPCA
alone and \sUnet{} for the hybrid U-Net on the same split (Table~\ref{tab:screens}, Fig.~\ref{fig:screens}). The
inner gradient steps initially fitted none of a scan's visible values; more steps and test-time refinement helped, but
the error then fell only slowly with network size, and the largest networks overfitted the deviations. As a hybrid,
functa added little to the PPCA prediction and nothing for regions or a missing family. Graph variants (Laplacian
coordinates, a learned vector per location, a graph-convolution decoder) did not help either. A learned vector per
location could in principle reproduce the linear PPCA, but fitting the latent by a few gradient steps did much worse
than PPCA's exact posterior.

\paragraph{Partial VAE.} The best partial VAE reached \sVaeHyb{} as a hybrid and \sVaeDev{} on the deviations, no
better than PPCA alone.

\paragraph{Age-dependent PPCA.} A population model per age, with kernel weights of log age, was worse than one model of
all ages at every bandwidth (\sAgePpca--\sAgePpcaMax{} against \sPpca{}): the deviations from the age-bin norm already
remove most of the age effect, and the local models lose scans.

\paragraph{Acquisition correction.} BCP scans of the two main acquisitions differed systematically in their deviations
(e.g.\ FA +0.42 SD and ODI $-0.27$ SD for \texttt{\bcpAcqBName} against \texttt{\bcpAcqAName}; software
documentation), but removing these offsets before imputation changed nothing (\rAcq{} vs \rDefault{}; paired
\cPRatio{}, Holm-adjusted $p=\cPPholm$). An offset that is constant within a scan is already carried by the scan's own
observed values and by the population prediction.

\paragraph{Context inputs and seed ensembles.} On seed 0 (54 scans), giving the transformer the distance to the nearest
observed value (blend \xDistBlend{} vs \rSeedZeroBlend; paired \xDistRatioBlend, $p=\xDistPBlend$) or the age and
acquisition (\xCtxBlend{}; paired \xCtxRatioBlend, $p=\xCtxPBlend$) did not help, nor did averaging three transformers
trained with different seeds (\xMembersBlend{}; paired \xMembersRatioBlend, $p=\xMembersPBlend$).

\subsection{A second cohort: pool the training of the networks, not the population model}\label{sec:pooling}

Pooling HBCD into the norm and the population model did not improve the BCP scans. In the 11-measure phase-3b
experiment (Appendix Table~\ref{tab:pooling}), pooling with the tensor and NODDI measures of HBCD left the error
unchanged (paired \vpTNRatio{}, CI \vpTNCi{}, $p=\vpTNP$), and pooling all measures including the free-water tensor
made it worse (\vpAllRatio{}, CI \vpAllCi{}), because HBCD's free-water fits differed from BCP's by more than a shift
and a scale; failed free-water fits with diffusivities far above that of free water had to be removed first. On the 7
measures, pooling the norm, PPCA and transformer made PPCA + SIREN worse (\rPooledDefault{} vs \rDefault{}; paired
\cQRatio{}), mostly for regions (\rPooledDefaultRegion) and a missing NODDI family (\rPooledDefaultWnoddi).

Training only the transformer on both cohorts, with the norm, PPCA and SIREN of BCP, helped: the transformer alone
improved from \rTfFour{} (400 epochs, BCP) to \rRefTf{} (paired \cRRatio{}, better on \cRBetter\%), and the
table blend from \rFinal{} to \rRefFinal{} (paired \cSRatio{}, better on \cSBetter\%). The HBCD-trained transformer was
not better for a missing NODDI family (\rRefFinalWnoddi{} vs \rFinalWnoddi{} for the blend).

\subsection{Error by measure and age}

\begin{table}[t]
\centering\footnotesize
\caption{Error relative to the age-bin mean per measure (four patterns, mean over age bins and seeds), and for a
whole missing family. $^{\ddagger}$First row: absolute RMSE of the age-bin mean (FA, NDI, ODI, FWF $\times10^{3}$;
MD, RD, AD $\times10^{6}$\,mm$^2$/s).}\label{tab:metric}
\setlength{\tabcolsep}{5pt}
\input{tables/tab_metric.tex}
\end{table}

\begin{figure}[t]
\centering
\includegraphics[width=\textwidth]{figures/fig_metric_heat.pdf}
\caption{Error relative to the age-bin mean per measure and pattern (mean over age bins and 3 seeds) for four
settings. ``Family'': the measure's family hidden (1.00 for phase 2, which cannot impute it).}\label{fig:metricheat}
\end{figure}

\paragraph{Measures.} The measures fell into two groups (Table~\ref{tab:metric}, Fig.~\ref{fig:metricheat}). The
diffusivities MD and RD and NDI vary smoothly in space and between neighbouring tracts and were imputed best (released
default \mMDFinal, \mRDFinal{} and \mNDIFinal{}); FA, AD, ODI and FWF depend more on the orientation and organization of
a tract and were imputed worst (FA \mFAFinal, AD \mADFinal, ODI \mODIFinal, FWF \mFWFFinal). For ODI and FWF in regions,
even the best methods remained near 0.8. For a missing family, the tensor measures were recovered well from NODDI
(released default: MD \mMDWmFinal, RD \mRDWmFinal, FA \mFAWmFinal), the NODDI measures less well from the tensor
(NDI \mNDIWmFinal, ODI \mODIWmFinal, FWF \mFWFWmFinal).

\paragraph{Age.} Relative to the age-bin mean, the error of the models varied little between age bins
(PPCA + SIREN \rDefaultBinA, \rDefaultBinB{} and \rDefaultBinC{} in 0--3, 4--9 and 10--60 months; fixed transformer
blend \rTfBlendFourBinA, \rTfBlendFourBinB{} and \rTfBlendFourBinC; Fig.~\ref{fig:age}a, b), and the ranking of the
methods was the same at every age. The per-age curves are noisy where few scans were validated (the youngest and
oldest ages). In absolute terms
(Table~\ref{tab:absolute}, Fig.~\ref{fig:age}c), the age-bin mean had its largest diffusivity errors in 0--3 months,
where values change fastest, and the released default reduced them most there.

\begin{figure}[t]
\centering
\includegraphics[width=\textwidth]{figures/fig_age.pdf}
\caption{Error by age. (a, b) Error relative to the age-bin mean per age in months (dots, area $\propto$ number of hidden
values; lines: Gaussian kernel smooth over log age) for PPCA alone, PPCA + SIREN, the transformer alone and the fixed
blend (400 epochs, 3 seeds). (c) Absolute RMSE of FA and MD for the age-bin mean and the released default per age
bin.}\label{fig:age}
\end{figure}

\begin{table}[t]
\centering\footnotesize
\caption{Absolute RMSE (four patterns, mean over patterns and seeds) per age bin; FA, NDI, ODI, FWF $\times10^{3}$, MD,
RD, AD $\times10^{6}$\,mm$^2$/s. Released default and HBCD-reference blend cross-fitted.}\label{tab:absolute}
\setlength{\tabcolsep}{4pt}
\input{tables/tab_absolute.tex}
\end{table}

\subsection{Uncertainty}

All model-based methods attach an SD to each imputed value. Calibrated per measure on held-out subjects, the SDs of
PPCA alone, the U-Net and the transformer were close to right for whole profiles and regions (root-mean-square of
error/SD \kPpcaWholeRms, \kUnetWholeRms{} and \kTfWholeRms; \kPpcaWholeIn\%, \kUnetWholeIn\% and \kTfWholeIn\% of
errors within 1.96 SD; Appendix Table~\ref{tab:calibration}), slightly too small for a missing family
(transformer \kTfFamilyRms). The SD of a value imputed by SIREN on the PPCA residuals is that of PPCA, an upper bound:
for gaps the errors were only \kDefGapRms{} of it on average. The SDs of the blend inherited this conservatism near
observed values (\kBlendGapRms{} for gaps) and were about right for regions (\kBlendRegionRms).

\subsection{Statistical summary}

Table~\ref{tab:paired} and Fig.~\ref{fig:paired} collect the pre-specified paired contrasts. All conclusions held
under the subject-level cluster bootstrap and the subject-level Wilcoxon test after Holm adjustment; the bootstrap CIs
were nearly identical to the normal-theory CIs, so the dependence between seeds and sessions inflated the precision of
the per-pair analysis only slightly. Differences of about 1\% (e.g.\ 40 or 120 vs 80 components) were detectable
because all settings saw the same hidden values; their practical relevance is small.

\begin{table}[t]
\centering\footnotesize
\caption{Paired per-scan comparisons over \nPairs{} scan $\times$ seed pairs (\nSubjPairs{} subjects): geometric mean of
the ratio of per-scan errors (compared / reference; below 1: compared is better), 95\% CI from the pairs and from a
subject-level cluster bootstrap, Holm-adjusted Wilcoxon $p$ over the \nContrasts{} contrasts (on the pairs / on subject
means), and the share of pairs where the compared setting is better.}\label{tab:paired}
\setlength{\tabcolsep}{3pt}
\resizebox{\textwidth}{!}{\input{tables/tab_paired.tex}}
\end{table}

\begin{figure}[t]
\centering
\includegraphics[width=0.8\textwidth]{figures/fig_paired.pdf}
\caption{Paired per-scan comparisons of Table~\ref{tab:paired}: geometric mean ratio (dot), 95\% CI over scan
$\times$ seed pairs (coloured line) and subject-level cluster bootstrap CI (grey bar).}\label{fig:paired}
\end{figure}

\section{Discussion}

\subsection{Principal findings}

The evaluation yields a simple picture. First, \emph{the scan alone does not identify a missing tract}: neighbouring
tracts differ systematically, and same-scan methods are only useful for short gaps along a profile. Second, \emph{the
normative reference is the main source of information}: imputing deviations from an age-specific norm, rather than
values, made all patterns imputable; modelling how deviations co-vary across all tracts and measures (PPCA) improved
whole profiles and regions further and made a whole missing measure family imputable. Third, \emph{local and global
information are complementary}: SIREN on the PPCA residuals within 4\,mm was the best near observed values, the
population model and amortized networks beyond. Fourth, \emph{amortized networks that read the observed values directly
add to the linear population model}, especially for a missing measure family, but only when they correct the PPCA
prediction and are blended with the local method; alone they were worse than PPCA + SIREN for the four patterns. The
transformer, with global attention over all tracts, was better than the graph U-Net with its local neighbourhoods.
Fifth, \emph{more data helps the networks but not the linear model}: training the transformer on a second, larger
cohort reduced the error by about 5\%, whereas pooling the second cohort into the norm and population model, after an
age-aware batch correction, made it worse.

\subsection{Interpretation}

Why does a linear population model of the deviations work so well, and why does functa fail where the transformer
succeeds? The PPCA posterior is the exact conditional expectation under a Gaussian model with one loading per cell;
with 80 components it captures the main modes of covariation of 37,716 cells from a few hundred scans. Functa and
the partial VAE replace this exact inference by a latent fitted with a few gradient steps or an encoder; with
hundreds of scans they did not learn a better prior than PPCA. The graph U-Net and the transformer, in contrast, read
each location's observed value directly and only need to learn a correction to the PPCA prediction; this is a much
easier learning problem, and their benefit was largest exactly where PPCA has least local information (beyond a few
millimetres and for a missing family). The blend weights quantify this division of labour: the transformer received
weight \wtCNearMin--\wtCNearMax{} next to observed values and \wtCFarMin--\wtCFarMax{} beyond 4\,mm.

The negative results are informative too. Correcting acquisition offsets did not help imputation because a
scan-wide offset is visible in the scan's observed values; this does not mean the offsets are irrelevant for group
analyses. Extra context (age, acquisition, distance) did not help the transformer because the observed values already
encode them. Averaging seeds did not help because the transformer's variance between seeds is small compared with
its bias relative to the truth on held-out subjects.

\subsection{Limitations}

\begin{itemize}
\item \emph{One normative cohort and internal validation.} All methods were developed and compared on BCP, with folds by
subject but with the hiding patterns, hyperparameters and defaults chosen on the same three seeds of validation
scans. The blend weights were cross-fitted for the reported scores, but the choice among blend variants was made on the
same data, so the final numbers may be slightly optimistic. An evaluation on an independent normative cohort is
needed \todo{e.g.\ HBCD scans as a test cohort, see experiments/requests.md}.
\item \emph{Synthetic missingness.} Values were hidden completely at random within each pattern. Real missingness is
not random: QC rejects tracts in scans with artifacts, more often in some tracts (Fig.~\ref{fig:data}b) and possibly in
some ages or subjects; the true values of rejected profiles may differ systematically from those of similar observed
profiles (missing not at random). The validation measures how well a method recovers plausible values, not whether it
corrects such biases.
\item \emph{Reference of the score.} Errors are relative to the age-bin mean, which favours methods on the age-bin norm
less in age bins where the norm is poor (0--3 months, where a continuous norm alone was 13\% better). The absolute RMSEs
(Table~\ref{tab:absolute}) complement the ratios.
\item \emph{Sample sizes.} The 0--3 month bin has only \bcpbinA{} BCP scans; 14 of them were validated per seed, and the
norm and models of that bin rest on few scans.
\item \emph{Dependence.} The 162 scan $\times$ seed pairs come from \nScansDistinct{} scans of \nSubjDistinct{}
subjects. We addressed this with a subject-level bootstrap and subject-level tests; the screens of extra inputs and
members used one seed only, and the functa/VAE screens one split.
\item \emph{Measures and atlas.} Results are for 7 measures of one atlas with dense tract axes; sparser tract sets give
fewer close neighbours, and the free-water tensor measures were left out after failed fits in HBCD.
\item \emph{Downstream impact.} We did not evaluate how imputation affects downstream statistical analyses (e.g.\ group
differences or age trajectories); imputed values are drawn towards the norm where the scan's own information is weak,
which can attenuate effects. Imputed values should be flagged (DTIPlayground writes masks and SDs) and sensitivity
analyses should exclude them.
\end{itemize}

\subsection{Practical recommendations}

\begin{enumerate}
\item \emph{Without a normative reference}, fill only short interior gaps (arc-linear along the profile, up to about
15\,mm); do not impute whole profiles or regions from the scan's own neighbours.
\item \emph{With a normative reference}, always impute deviations from an age-appropriate norm rather than values.
\item \emph{With a normative cohort of a few hundred scans}, fit a population model (PPCA, about 80 components) of the
deviations over all tracts and measures, impute its residuals locally within about 4\,mm (SIREN, or $k$NN without a
GPU at about 8\% more error), and use the PPCA prediction beyond. This needs no network training and was within 8\% of
the best setting.
\item \emph{If a GPU is available}, blend this with a transformer trained to correct the PPCA prediction, with weights
increasing with the distance to the nearest observed value (the default of DTIPlayground 0.9.8). This is most valuable
when whole measure families are missing.
\item \emph{If a larger external cohort on the same tracts exists}, use it to train the transformer (after a batch
correction), but keep the norm and population model of the target cohort.
\item \emph{Report} the share of imputed values per tract and measure, keep the imputation masks and SDs, and check
that conclusions hold without imputed values. Do not rely on imputed values for group differences between normative
and clinical scans, since imputed values are drawn towards the norm.
\end{enumerate}

\section{Conclusion}

In a developmental cohort with dense atlas tracts, missing diffusion profile values are best imputed as deviations
from an age-specific norm, using a joint population model of all tracts and measures, refined locally near observed
values and blended with an amortized transformer beyond them. This reduced the error to \rFinal{} of the age-bin
mean's overall (\rRefFinal{} with a second cohort for training), and to \rFinalWdti{} and \rFinalWnoddi{} when a whole
tensor or NODDI fit is missing. Several more complex alternatives did not help. All methods, the validation protocol
and the defaults are available in DTIPlayground.

\section*{Data and code availability}

DTIPlayground (including \texttt{dmrifiberprofile impute-profiles} and its \texttt{--validate} protocol) is open source
at \url{https://github.com/NIRALUser/DTIPlayground}, on PyPI (\texttt{dtiplayground}, version 0.9.8) and on Docker Hub
\todo{Docker image name}. The validation of every phase is documented in the package documentation
(\texttt{docs/ImputeValidation.rst}). The scripts that compute every number, table and figure of this paper from the
validation result files are in \texttt{papers/ProfileImputation/scripts} of the repository \todo{confirm that the paper
folder will be public}. BCP data are available through \todo{NDA collection / BCP data access}; HBCD data through
\todo{HBCD data release / NBDC}. \todo{Whether derived profile tables can be shared.}

\section*{Acknowledgements}
\todo{Funding (BCP, HBCD, grants), data use acknowledgements required by BCP and HBCD.}

\section*{Author contributions}
\todo{CRediT statement.}

\section*{Declaration of competing interests}
\todo{Competing interests.}

\bibliographystyle{plainnat}
\bibliography{refs}

\clearpage
\appendix
\section{Supplementary material}
\setcounter{figure}{0}\renewcommand{\thefigure}{A\arabic{figure}}
\setcounter{table}{0}\renewcommand{\thetable}{A\arabic{table}}

\begin{figure}[h]
\centering
\includegraphics[width=\textwidth]{figures/fig_siren_grid.pdf}
\caption{SIREN on the deviations from the age-bin norm: error relative to the age-bin mean against the training epochs,
per input (rows), $\omega_0$ (colours) and pattern (columns). Earlier protocol: 30 scans (10 per age bin), one seed, 11
measures (file \texttt{docs/\_static/impute\_validation/phase2/stratified\_overall.csv} of the
repository).}\label{fig:sirengrid}
\end{figure}

\begin{figure}[h]
\centering
\includegraphics[width=\textwidth]{figures/fig_cap.pdf}
\caption{Distance limit of the local methods: error relative to the age-bin mean when the local method fills cells up to
the limit and the PPCA prediction (or, for $k$NN on the deviations, the age-bin mean) beyond. Earlier protocol: 30 scans,
one seed, 11 measures (\texttt{docs/\_static/impute\_validation/phase3/grid\_caps\_long.csv}).}\label{fig:cap}
\end{figure}

\begin{table}[h]
\centering\footnotesize
\caption{Pooling HBCD and continuous age norms (phase-3b protocol: 3 seeds $\times$ 54 BCP scans, 11 measures including
the free-water tensor measures; BCP scans scored in BCP units against the BCP-alone age-bin mean). ``Pooled
tensor+NODDI'': HBCD's free-water tensor measures left out. All settings with norm except ``Age norm alone'' are PPCA
+ SIREN.}\label{tab:pooling}
\input{tables/tab_pooling.tex}
\end{table}

\begin{table}[h]
\centering\footnotesize
\caption{Calibration of the predicted SDs: root-mean-square of error/SD (1: right; below 1: SDs too large) and share of
errors within 1.96 SD (95\%: right), mean over measures, 3 seeds.}\label{tab:calibration}
\input{tables/tab_calibration.tex}
\end{table}

\begin{table}[h]
\centering\footnotesize
\caption{Blend weight tables (weight $w$ of the transformer) refitted from the validation dumps; identical to the
tables shipped in DTIPlayground 0.9.8 (\texttt{ENSEMBLE\_TABLES}).}\label{tab:weights}
\setlength{\tabcolsep}{3pt}
\resizebox{\textwidth}{!}{\input{tables/tab_weights.tex}}
\end{table}

\end{document}
